\section{结论}\label{sec:conclusion}
针对复杂场景下多模态情感识别的三个问题，本文形成了由原始媒体特征构建、局部缺失预测到层级证据解释的完整方案。问题一对附件1三模态媒体分别构造带时间支持的原生特征，以最大时间重叠映射至50个公共时间窗，并保留有效掩码和来源索引，得到可追溯的时序表示。其辅助交叉验证帮助比较构建方案，但不作为独立泛化结论。

问题二直接利用附件2预计算对齐特征，建立三分类与连续强度的双任务基线，并用轻量时间注意力残差池化补充均值池化的时间选择能力。在验证集的完整输入和54种连续缺失条件下，三次初始化的平均综合鲁棒分数取得小幅改进；宏平均F1、MAE与相关系数的平均方向较好，准确率略有下降，而且收益随初始化与缺失场景变化。该结果支持选择低参数开销的模型，同时要求对具体应用场景继续检查误差。

问题三提出层级式多模态证据归因框架，通过八种模态联盟的精确贡献、模态间交互、连续时间遮挡和删除干预，构建由模态到局部区间的解释。附件4的20条无标签样本均生成预测与解释卡；文本证据可映射到经核验的原始字符片段，语音与视觉目前仅能定位到特征槽位。后续工作应以独立来源的音视频时间映射和有标签外部评价扩展现有证据边界，而不能把无标签预测或模型干预响应当作真实情感的正确性与因果证明。
